\section{Limitations and Potential Solutions}
\label{supp:limit}

While \molmoactt is all quite capable as a general-purpose action reasoning model, it is not without limitations. In the following sections, we discuss some of these limitations and potential solutions.

\paragraph{Action chunks and the absence of real-time re-chunking.}
\molmoactt predicts fixed-horizon action chunks (30 steps at 30~Hz for YAM/SO-100/101, 15 at 15~Hz for DROID, 10 at 10~Hz for LIBERO) and executes them open-loop before re-querying the policy. This has two consequences worth acknowledging. First, motion smoothness across chunk boundaries is not enforced anywhere in the pipeline: the flow-matching expert is trained to denoise each chunk independently, with no continuity loss tying the end of chunk $t$ to the start of chunk $t+1$. In practice this can produce visible velocity or acceleration discontinuities at the seams, particularly when the VLM context shifts between queries (new observation, slightly different attention state). Second, because the policy commits to a full chunk before observing its consequences, it cannot react within-chunk to perturbations, contact events, or its own tracking error --- the 55.79~Hz number in Sec \ref{speed} is amortized chunk throughput, not closed-loop reactivity.

\paragraph{Embodiment-specific zero-shot deployment.}
The zero-shot deployment story for \textsc{MolmoAct2} is currently tied to the three embodiments for which we have large-scale training data: \molmoacttyam, \molmoacttso, and \molmoacttdroid can each be deployed out-of-the-box on their respective robot setups (bimanual YAM, SO-100/101, and DROID Franka) Sec. \ref{sec:exp-outbox}. This is a direct consequence of our data strategy --- concentrating collection and curation effort on these three platforms to span the low-to-medium cost range --- but it also bounds what zero-shot means here. \textsc{MolmoAct2} is not a universal controller that transfers without adaptation to arbitrary new embodiments: deployment on robots outside this set (e.g., other bimanual platforms, dexterous hands, mobile manipulators with different kinematics, or humanoids) requires effective fine-tuning on target-embodiment demonstrations, as shown in our fine-tuning studies in Sec \ref{sec:exp-ft}. Extending the zero-shot envelope to a broader range of embodiments is a matter of scaling data collection along the same axes we have begun here, and we view the released datasets and training recipe as a foundation for that effort rather than a finished product.

